\section{Introduction}
Lipases catalyse esterifications, transesterifications and hydrolyses that are central to the valorisation of biomass-derived lipids \cite{gamayurova2021,sousa2021,paiva2000}. Many of these processes are complex in the kinetic sense: several substrates compete for the enzyme, or the product of one reaction is the substrate of the next. The esterification of trimethylolpropane (TMP) with fatty acids, a key step towards fully bio-based biolubricants because both TMP and the fatty acids can be produced from biomass \cite{sayed2025,lennen2013}, is a prototypical case. TMP (Z) is converted to the monoester (M), the diester (D) and the triester (T) in three consecutive reactions, each of which follows the Ping Pong bi bi mechanism \cite{cleland1963,cornishbowden2012}.

Reactor design and optimisation require kinetic models of such processes, but fully mechanistic Ping Pong rate equations contain many correlated constants that cannot be estimated reliably from concentration--time profiles \cite{alhaque2012,bornadel2013}. Nishimura et al. \cite{nishimura2026} recently proposed in \emph{Biomass} a two-step semi-mechanistic approach that circumvents this problem. In the first step the selectivities of the enzyme for the second and third esterification ($S_M$, $S_D$) are estimated with an ``irreversible fingerprinting model'' in which the degree of reaction $c$ replaces time as the independent variable, so that enzyme saturation, inhibition and denaturation cancel out of the equations \cite{mitchell2022a,mitchell2022b,mitchell2023}. In the second step $S_M$ and $S_D$ are fixed and a dynamic model with lumped empirical constants ($k_A$, $F_A$, $F_{AZ}$, $F_{AM}$, $F_{AD}$) is fitted against time. Four literature case studies \cite{akerman2011,bornadel2013,tao2015,mao2024} were described ``well'', judged by the minimised sum of squares $F_\mathrm{obj}$ and the Bayesian information criterion (BIC).

The approach is attractive because it is transparent and because every data point is published in the Supplementary Material. Complete public data are rare in kinetic modelling and offer an opportunity that we take up here: the method can be audited, reproduced to the last printed digit and improved on exactly the same data. Three features of the original analysis motivate this study, and they are connected. (i) Pseudo-processivity, the re-acylation of a freshly formed ester before it leaves the active site, which mass-transfer limitation produces in the data of Bornadel et al. \cite{bornadel2013,mitchell2023}, was deliberately excluded, so the selectivities of that case are, in the authors' own words, only ``apparent''. (ii) Enzyme deactivation was handled by a separate rate constant for each reuse batch, never within a batch. (iii) The models were judged only in-sample, on the data used for fitting, and because $S_M$ and $S_D$ are treated as known in step 2, the two parameters estimated from the same data in step 1 are absent from the reported BIC. The connection is this: the two-step method buys its stability by fixing the selectivities from the $c$-domain; if the $c$-domain model omits a mechanism, the resulting error is locked into step 2, and an in-sample criterion applied to step 2 alone cannot reveal it.

This observation defines the main line of the paper, which we state as a single argument and then test step by step (Table~\ref{tab:roadmap}). First, the published analysis is reproduced from the public data and audited, which shows what in-sample statistics can and cannot establish. Second, the misfit left by the two-step model is examined for its mechanistic signature. Third, we introduce \textbf{PPK-FIP} (\emph{pseudo-processive kinetics with a fingerprint-informed prior}), which keeps the Ping Pong core and the central idea of Nishimura et al., that selectivities are best learned against the degree of reaction, but (a) adds pseudo-processivity and first-order deactivation as explicit, testable mechanisms, (b) replaces the hard ``fix-then-fit'' constraint by a calibrated Bayesian prior centred on the fingerprint estimate, (c) averages the simplified and full lumped-constant forms with BIC weights, choosing the prior hyper-parameters by leave-one-case-out meta-validation, never with the target data, and (d) exploits an exact reparameterisation of the shared Ping Pong denominator that makes global optimisation and exhaustive cross-validation affordable. Fourth, PPK-FIP is compared with the original method on the paper's own criteria and, decisively, out-of-sample, and an ablation attributes the gain to its components. Fifth, cross-condition transfer and Monte-Carlo experiments with known truth test whether the added mechanisms are genuine properties of the reaction systems rather than extra flexibility. Sixth, the mechanistic estimators are set against modern black-box learners (Gaussian processes, MLP ensembles, the pretrained tabular foundation model TabICL \cite{qu2025tabicl} and neural ordinary differential equations \cite{chen2018}) and then combined with them in a hybrid, PPK-FIP-GP. Every figure and table of the paper belongs to one of these six steps, and after each result we state what it contributes to the argument.

\begin{table}[!htbp]
\centering\small
\caption{Main line of the study: six questions, the experiments that answer them, and where the evidence is shown. Figures~\ref{fig:workflow}, \ref{fig:data} and \ref{fig:bench} (Section~\ref{sec:methods}) describe the method and the data that all six steps share.}\label{tab:roadmap}
\begin{adjustbox}{max width=\linewidth}
\begin{tabular}{@{}c p{5.6cm} p{4.6cm} p{4.4cm}@{}}
\toprule
Step & Question & Experiment & Evidence \\
\midrule
1 & Is the published analysis reproducible from its public data, and is in-sample fit a reliable guide to the model? & Data audit; re-fit of both steps; multistart and profile-likelihood analysis & Tables~\ref{tab:data}--\ref{tab:reprodyn}; Figures~\ref{fig:repro}, \ref{fig:landscape} \\
2 & Which mechanisms does the remaining misfit point to? & Generalised model fitted to the full data; residual and composition-path diagnostics & Tables~\ref{tab:insample}, \ref{tab:params}; Figures~\ref{fig:fits}--\ref{fig:resid} \\
3 & Does the generalised model predict unseen data, not merely describe seen data? & Leave-one-time-point-out cross-validation; forward extrapolation & Tables~\ref{tab:loto}, \ref{tab:fcst}; Figures~\ref{fig:loto}--\ref{fig:fcst} \\
4 & Which component carries the gain, and must the fingerprint be used as a prior? & Component ablation; prior width; multi-task pooling & Tables~\ref{tab:ablation}, \ref{tab:pooling}; Figures~\ref{fig:ablation}, \ref{fig:pooling} \\
5 & Are the added mechanisms genuine properties of the reaction systems? & Cross-condition transfer; Monte-Carlo truth recovery & Tables~\ref{tab:transfer}, \ref{tab:sim}; Figures~\ref{fig:transfer}, \ref{fig:sim} \\
6 & Is mechanistic structure still needed, and how should machine learning be used? & Black-box baselines; Gaussian-process discrepancy hybrid & Tables~\ref{tab:lotobb}--\ref{tab:hybrid}; Figures~\ref{fig:blackbox}, \ref{fig:hybrid} \\
\bottomrule
\end{tabular}
\end{adjustbox}
\end{table}
